% =====================================================================
% POLI 450A (2026) --- shared TEMPLATE layer.
%
% Input from each lecture near the top of its file, BEFORE
% ../course_theme.tex (metropolis must be selected before the course
% palette overrides it):
%
%   \documentclass[xcolor=svgnames,9pt,english,aspectratio=169]{beamer}
%   \newcommand{\handout}{0}
%   \input{../course_preamble.tex}
%   \input{../course_theme.tex}
%   \graphicspath{{./}{./figs/}}
%   \title[Short title]{450A Political Methodology I}
%   \subtitle{Lecture N: ...}
%
% Modeled on the 2026 Northwestern workshop template
% (Teaching/PanelData/_latex/2026/202608_Northwestern/workshop_preamble.tex).
%
% SCOPE: theme, beamer layout, packages, the roadmap hooks, and --- unlike
% the Northwestern decks --- the notation macros. The 450A lectures
% grew from one common ancestor, so their macro blocks agree wherever they
% overlap; merging them is safe. Three points where they had drifted are
% resolved here, and each is flagged inline below.
% =====================================================================

\usetheme{metropolis}
\metroset{sectionpage=none, subsectionpage=none, numbering=none}
\setbeamertemplate{navigation symbols}{}

\setbeamersize{text margin left=5mm, text margin right=5mm}
\setbeamertemplate{itemize/enumerate body begin}{\leftmargini=1.5em\relax}
\setbeamertemplate{itemize/enumerate subbody begin}{\leftmarginii=1.5em\relax}
\setbeamertemplate{itemize/enumerate subsubbody begin}{\leftmarginiii=1.5em\relax}

% ----------------------------------------------------------- top alignment
% Every deck passes the t class option, so beamer applies the frame key
% below to every frame (explicit [plain,c] frames — title page, roadmaps —
% stay centered). Beamer's stock t key puts a rigid .2cm pad above the
% content; frames sized to this template's full-height ceilings then run
% Overfull \vbox by up to that pad. Same key as beamerbaseframe.sty, with
% one change: the pad may shrink to zero when the frame is full.
\makeatletter
\define@key{beamerframe}{t}[true]{% top
  \beamer@frametopskip=.2cm plus .5\paperheight minus .2cm\relax%
  \beamer@framebottomskip=0pt plus 1fill\relax%
  \beamer@frametopskipautobreak=\beamer@frametopskip\relax%
  \beamer@framebottomskipautobreak=\beamer@framebottomskip\relax%
}

% ------------------------------------------------- slide-top check (2026-10-11)
% A slide whose content is too tall for the frame is squeezed: the pad above
% shrinks and its content starts higher than on other slides, so a figure
% "wobbles" when the class steps from one slide to the next (Lecture 4-1,
% slides 1-2, 2026-10-11). The marker below records where each slide's
% content starts (a whatsit: it adds no space and moves nothing).
% `python3 ../check_slides.py lec<N>` then lists every slide that starts
% higher than the deck's normal position. That list must be empty.
% Recorded per slide: page, frame number, content-top y (sp), the room left
% below the title bar, and P for a plain frame (title page, roadmaps).
\providecommand\slidetop[5]{}
\AtBeginDocument{\immediate\write\@auxout{\string\providecommand\string\slidetop[5]{}}}
\def\beamer@markframetop{%
  \pdfsavepos
  \edef\beamer@marktmp{\noexpand\write\noexpand\@auxout{\string\slidetop
    {\noexpand\thepage}{\noexpand\insertframenumber}{\noexpand\the\noexpand\pdflastypos}%
    {\number\beamer@frametextheight}{\ifbeamer@plainframe P\else N\fi}}}%
  \beamer@marktmp}
\RequirePackage{etoolbox}
\patchcmd{\endbeamer@frameslide}
  {\vskip\beamer@frametopskip}
  {\vskip\beamer@frametopskip\beamer@markframetop}
  {}{\PackageWarning{course_preamble}{slide-top marker not installed}}
\makeatother

% ---------------------------------------------------------------- packages
\usepackage{xcolor}
\usepackage{colortbl}
\usepackage{amsmath}
\usepackage{amsfonts}
\usepackage{amssymb}
\usepackage{bm}
\usepackage{graphicx}
\usepackage{multirow}
\usepackage{fancyvrb}
\usepackage{listings}
\usepackage{tikz}
\usetikzlibrary{arrows.meta,positioning,decorations.pathreplacing}

% --------------------------------------------------------- figure fitting
% Default bounding box for every \includegraphics that does not override it.
% With keepaspectratio the image is scaled to fit *inside* width x height, so
% a figure can never run past the side margins or off the bottom of a slide.
% Options given at the call site override these on a per-key basis: e.g.
% [width=.6\linewidth] narrows the box but keeps the height ceiling.
% \figmaxheight is set at \begin{document}: beamer only computes the final
% \textheight once the head/foot templates are known. 0.80 is just under what a
% frame with a title and one line of text can hold (measured: 0.834\textheight);
% a frame that needs more room for text can lower it locally.
\newlength{\figmaxheight}
\AtBeginDocument{\setlength{\figmaxheight}{0.80\textheight}}
\setkeys{Gin}{width=\linewidth,height=\figmaxheight,keepaspectratio}

% -------------------------------------------------------- roadmap frames
% Identical in all lectures. A lecture that does not want the
% subsection roadmap clears it after the \input with \AtBeginSubsection[]{}.
% hideothersubsections is load-bearing, not cosmetic: lecture 4 had 13
% subsections before it was split into 4-1, 4-2 and 4-3 (2026-10-09), and a
% full contents list ran off the top and bottom of the slide. Hiding the other sections' subsections caps the list at 10 lines.
\subject{Talks}
% No back-to-back roadmaps (Yiqing, 2026-10-11): when a subsection starts right
% after its section's roadmap, with no content slide in between, its own
% roadmap would repeat the one just shown (5-1's handout showed "Two Cases"
% twice), so it is skipped. The section roadmap records the frame count; the
% subsection roadmap appears only if a content slide has come since.
\makeatletter
\def\course@secframe{-1}
\AtBeginSection[]
{
  \begin{frame}[plain,c]
    \frametitle{}
    \tableofcontents[currentsection,hideothersubsections]
    \addtocounter{framenumber}{-1}
  \end{frame}
  \xdef\course@secframe{\the\c@framenumber}
}
\AtBeginSubsection[]
{
  \ifnum\c@framenumber=\course@secframe\relax\else
  \begin{frame}[plain,c]
    \frametitle{}
    \tableofcontents[currentsection,currentsubsection,hideothersubsections]
    \addtocounter{framenumber}{-1}
  \end{frame}
  \fi
}
\makeatother

% ------------------------------------------------------ theorem environments
\newtheorem{assumption}{Assumption}
\newtheorem{iass}{Identification Assumption}
\newtheorem{ires}{Identification Result}
\newtheorem{estm}{Estimand}
\newtheorem{esti}{Estimator}
\newtheorem{question}{Question}
\newtheorem{answer}{Answer}
\newtheorem{proposition}{Proposition}

% =====================================================================
% Notation
% =====================================================================
\definecolor{dgreen}{rgb}{0,.6,0.8}

\renewcommand\r{\right}
\renewcommand\l{\left}
\newcommand{\bs}{\boldsymbol}
\newcommand{\mb}{\mathbf}
\newcommand{\real}{\mathbb{R}}
\newcommand{\cN}{\mathcal{N}}
\newcommand{\E}{\mathbb{E}}
\newcommand{\V}{\mathbb{V}}

% MERGE NOTE 1. Lecture 2-4 set \epsilon to \varepsilon; the other twelve
% left it at the upright default. Adopting \varepsilon course-wide so the
% error term looks the same in every lecture.
\renewcommand{\epsilon}{\varepsilon}

% MERGE NOTE 2. Lecture 2-4 wrote its vectors as \bm (bold italic); the other
% twelve use \mathbf (bold upright). \mathbf wins --- it is the convention in
% eleven of the thirteen decks and in the course's matrix lectures. Only
% \iid, \E, \cN and \Var from that block are actually used in 2-4's body.
\renewcommand{\u}{\ensuremath{\mb{u}}}
\renewcommand{\v}{\ensuremath{\mb{v}}}
\newcommand{\U}{\ensuremath{\mb{U}}}
\newcommand{\M}{\ensuremath{\mb{M}}}
\newcommand{\X}{\ensuremath{\mb{X}}}
\newcommand{\x}{\ensuremath{\mb{x}}}
\newcommand{\y}{\ensuremath{\mb{y}}}
\newcommand{\Z}{\ensuremath{\mb{Z}}}
\newcommand{\C}{\ensuremath{\mb{C}}}
\newcommand{\D}{\ensuremath{\mb{D}}}
\newcommand{\I}{\ensuremath{\mb{I}}}
\newcommand{\W}{\ensuremath{\mb{W}}}
\renewcommand{\a}{\ensuremath{\mb{a}}}
\renewcommand{\P}{\ensuremath{\mb{P}}}
\newcommand{\one}{\ensuremath{\mb{1}}}
\newcommand{\zero}{\ensuremath{\mb{0}}}

\renewcommand{\b}{\ensuremath{\bs{\beta}}}
\newcommand{\be}{\ensuremath{\bs{\beta}}}
\newcommand{\e}{\ensuremath{\bs{\epsilon}}}
\newcommand{\eps}{\ensuremath{\bs{\epsilon}}}
\newcommand{\eeta}{\ensuremath{\bs{\eta}}}
\newcommand{\bhat}{\ensuremath{\hat{\bs{\beta}}}}
\newcommand{\ehat}{\ensuremath{\hat{\e}}}
\newcommand{\Sig}{\ensuremath{\bs{\Sigma}}}
\newcommand{\SigInv}{\ensuremath{\bs{\Sigma^{-1}}}}

\newcommand{\XX}{\ensuremath{\X'\X}}
\newcommand{\XXinv}{\ensuremath{\left(\XX\right)^{-1}}}
\newcommand{\hatmat}{\ensuremath{\X \XXinv \X'}}
\newcommand{\hatsig}{\ensuremath{\hat{\sigma}^2}}
\newcommand{\sig}{\ensuremath{\sigma^2}}

\newcommand{\cA}{\mathcal{A}}
\newcommand{\cJ}{\mathcal{J}}
\newcommand{\cL}{\mathcal{L}}
\newcommand{\cLL}{\mathcal{LL}}
\newcommand{\cll}{\ell}
\newcommand{\cQ}{\mathcal{Q}}
\newcommand{\cP}{\mathcal{P}}

% MERGE NOTE 3. \var / \Var and \cov / \Cov were spelled with \text in some
% lectures and \mbox in others. Same output; \text kept.
\newcommand{\var}{\text{Var}}
\newcommand{\Var}{\text{Var}}
\newcommand{\cov}{\text{Cov}}
\newcommand{\Cov}{\text{Cov}}
\newcommand{\cor}{\text{Cor}}
\newcommand{\tr}{\text{tr}}
\newcommand{\avepop}{\frac{1}{n}\sum_{i=1}^n}

\DeclareMathOperator*\plim{plim}
\DeclareMathOperator\rank{rank}
\DeclareMathOperator*\argmax{argmax}
\DeclareMathOperator*\argmin{argmin}
\DeclareMathOperator{\sgn}{sgn}
\newcommand{\diag}{\mathop{\mathrm{diag}}}

\newcommand\ud{\mathrm{d}}
\newcommand\dist{\buildrel\rm d\over\sim}
\newcommand\ind{\stackrel{\rm indep.}{\sim}}
\newcommand\iid{\stackrel{\rm i.i.d.}{\sim}}
\newcommand\logit{{\rm logit}}

\newcommand\indep{\protect\mathpalette{\protect\independenT}{\perp}}
\def\independenT#1#2{\mathrel{\rlap{$#1#2$}\mkern2mu{#1#2}}}

% ------------------------------------------------------------- text helpers
\newcommand{\bi}{\begin{itemize}}
\newcommand{\ei}{\end{itemize}}
\newcommand{\noi}{\noindent}
\newcommand{\red}{\textcolor{Cardinal}}
\newcommand{\blue}{\textcolor{CourseBlue}}
\newcommand{\gray}{\textcolor{gray}}
\newcommand\spacingset[1]{\renewcommand{\baselinestretch}{#1}\small\normalsize}

\hyphenation{tyamamot}


% ------------------------------------------------ figure-position check (2026-10-11)
% Records where every \includegraphics lands and its size (a whatsit; it moves
% nothing). check_slides.py flags two consecutive slides that show a figure of
% the same size in different places: the class sees it jump (Lecture 3, slides
% 18 to 19: one frame had \vspace{.3cm} above the figure, the next did not),
% and two consecutive slides that show the same file at different sizes or
% places (Lecture 6-3, slides 10 to 11, the ballot image). The file name is the
% seventh field.
\makeatletter
\providecommand\figpos[7]{}
\AtBeginDocument{\immediate\write\@auxout{\string\providecommand\string\figpos[7]{}}}
\newbox\course@figbox
\let\course@origGinclude\Ginclude@graphics
\def\Ginclude@graphics#1{%
  \setbox\course@figbox\hbox{\course@origGinclude{#1}}%
  \pdfsavepos
  \edef\course@figtmp{\noexpand\write\noexpand\@auxout{\string\figpos
    {\noexpand\thepage}{\noexpand\insertframenumber}%
    {\noexpand\the\noexpand\pdflastxpos}{\noexpand\the\noexpand\pdflastypos}%
    {\number\wd\course@figbox}{\number\ht\course@figbox}{\detokenize{#1}}}}%
  \course@figtmp
  \box\course@figbox}
\makeatother
